% =============================================================================
% Présentation --- Détection de fraude par carte bancaire (Machine Learning)
% Version améliorée : images correctement dimensionnées, UI professionnelle
%
% Compilation depuis la RACINE du projet (pdflatex x2) — figures = dossier outputs/ :
%   pdflatex -output-directory=latex_build latex/presentation_fraude_complete.tex
%   pdflatex -output-directory=latex_build latex/presentation_fraude_complete.tex
% Prérequis : exécuter 01_EDA.ipynb et 02_Modeling.ipynb pour remplir outputs/
% =============================================================================
\documentclass[aspectratio=169,11pt]{beamer}

% ─────────────────── Thème & couleurs ────────────────────────────────────────
\usetheme{AnnArbor}
\usecolortheme{default}
\usefonttheme{professionalfonts}
\setbeamertemplate{navigation symbols}{}

% Palette principale
\definecolor{MLDeep}   {RGB}{13,  47,  90}   % bleu marine profond
\definecolor{MLPrimary}{RGB}{21,  67, 115}   % bleu institutionnel
\definecolor{MLLight}  {RGB}{235,243,253}    % bleu très clair (fond)
\definecolor{MLAccent} {RGB}{183, 28,  28}   % rouge alerte
\definecolor{MLMuted}  {RGB}{96, 110, 125}   % gris bleuté
\definecolor{MLGreen}  {RGB}{ 27,120,  80}   % vert succès
\definecolor{MLAmber}  {RGB}{195,125,  10}   % ambre avertissement

% Couleurs Beamer
\setbeamercolor{structure}           {fg=MLPrimary}
\setbeamercolor{frametitle}          {bg=MLDeep,   fg=white}
\setbeamercolor{frametitle right}    {bg=MLPrimary}
\setbeamercolor{title}               {fg=MLDeep}
\setbeamercolor{subtitle}            {fg=MLMuted}
\setbeamercolor{author in head/foot} {bg=MLPrimary, fg=white}
\setbeamercolor{title in head/foot}  {bg=MLDeep,    fg=white!85}
\setbeamercolor{date in head/foot}   {bg=MLDeep,    fg=white!70}
\setbeamercolor{section in head/foot}{bg=MLLight,   fg=MLDeep}
\setbeamercolor{section in toc}      {fg=MLPrimary}
\setbeamercolor{subsection in toc}   {fg=black!70}
\setbeamercolor{itemize item}        {fg=MLPrimary}
\setbeamercolor{itemize subitem}     {fg=MLMuted}
\setbeamercolor{enumerate item}      {fg=MLPrimary}
\setbeamercolor{block title}         {bg=MLDeep,         fg=white}
\setbeamercolor{block body}          {bg=MLLight,        fg=black}
\setbeamercolor{block title alerted} {bg=MLAccent!85,    fg=white}
\setbeamercolor{block body alerted}  {bg=MLAccent!8,     fg=black}
\setbeamercolor{block title example} {bg=MLGreen!80!black,fg=white}
\setbeamercolor{block body example}  {bg=MLGreen!7,      fg=black}

% Polices
\setbeamerfont{title}     {size=\LARGE, series=\bfseries}
\setbeamerfont{subtitle}  {size=\normalsize, shape=\itshape}
\setbeamerfont{frametitle}{size=\large,  series=\bfseries}
\setbeamerfont{block title}{size=\small, series=\bfseries}

% Items
\setbeamertemplate{itemize items}[circle]
\setbeamertemplate{itemize subitem}[triangle]
\setbeamertemplate{enumerate items}[default]
\setbeamertemplate{section in toc}[sections numbered]
\setbeamertemplate{caption}[numbered]

% Pied de page compact
\setbeamertemplate{footline}{%
  \leavevmode%
  \hbox{%
    \begin{beamercolorbox}[wd=.30\paperwidth,ht=2.5ex,dp=1ex,left,leftskip=1em]{author in head/foot}%
      \usebeamerfont{author in head/foot}\insertshortauthor
    \end{beamercolorbox}%
    \begin{beamercolorbox}[wd=.50\paperwidth,ht=2.5ex,dp=1ex,center]{title in head/foot}%
      \usebeamerfont{title in head/foot}\insertshorttitle
    \end{beamercolorbox}%
    \begin{beamercolorbox}[wd=.20\paperwidth,ht=2.5ex,dp=1ex,right,rightskip=1em]{date in head/foot}%
      \usebeamerfont{date in head/foot}%
      \insertframenumber\,/\,\inserttotalframenumber
    \end{beamercolorbox}%
  }%
  \vskip0pt%
}

% ─────────────────── Paquets ─────────────────────────────────────────────────
\usepackage[T1]{fontenc}
\usepackage[french]{babel}
\usepackage{microtype}
\usepackage{graphicx}
\usepackage{booktabs}
\usepackage{array}
\usepackage{colortbl}
\usepackage{xcolor}
\usepackage{tikz}
\usepackage{hyperref}

\hypersetup{
  colorlinks=true,
  linkcolor=MLPrimary,
  urlcolor=MLAccent,
  citecolor=MLMuted,
}

% ─────────────────── Chemin des figures (toutes dans outputs/) ─────────────
\newcommand{\figdir}{outputs/}  % pdflatex depuis la racine du projet
% \renewcommand{\figdir}{../outputs/}  % si compilation depuis latex/

% ─────────────────── Macro d'inclusion sécurisée ─────────────────────────────
% Usage : \safeimg[width_ratio]{filename}  (défaut 0.85)
% Calcule aussi une hauteur max pour ne jamais déborder du frame.
\newlength{\imgmaxw}
\newlength{\imgmaxh}
\newcommand{\safeimg}[2][0.85]{%
  \setlength{\imgmaxw}{#1\linewidth}%
  \setlength{\imgmaxh}{0.58\textheight}%  garde place pour titre + pied
  \IfFileExists{\figdir#2}{%
    \begin{center}%
      \begin{tikzpicture}%
        \node[%
          inner sep=4pt,%
          rounded corners=3pt,%
          draw=MLPrimary!30,%
          line width=0.5pt,%
          fill=white,%
        ] {%
          \includegraphics[%
            width=\imgmaxw,%
            height=\imgmaxh,%
            keepaspectratio,%
          ]{\figdir#2}%
        };%
      \end{tikzpicture}%
    \end{center}%
  }{%
    \begin{center}%
      \begin{tikzpicture}
        \node[%
          draw=MLAccent!50,%
          fill=MLAccent!6,%
          rounded corners=4pt,%
          inner sep=10pt,%
          text width=0.80\linewidth,%
          align=center,%
          font=\small,%
        ] {%
          \textcolor{MLAccent}{\textbf{\small Figure manquante}}\\\smallskip
          \texttt{\small\figdir\detokenize{#2}}\\[4pt]
          \textcolor{MLMuted}{Générer \texttt{outputs/} via \texttt{01\_EDA.ipynb} / \texttt{02\_Modeling.ipynb}}%
        };
      \end{tikzpicture}%
    \end{center}%
  }%
}

% Macro pour deux images côte à côte dans des colonnes (évite tout débordement)
\newcommand{\twoimgs}[2]{%
  \begin{columns}[T,onlytextwidth]
    \column{0.5\textwidth}%
      \safeimg[0.96]{#1}%
    \column{0.5\textwidth}%
      \safeimg[0.96]{#2}%
  \end{columns}%
}

% ─────────────────── Métadonnées ─────────────────────────────────────────────
\title[Détection de fraude ML]{Détection de fraude par carte bancaire}
\subtitle{Projet Machine Learning --- Polytech Semestre~4}
\author[Jegham \and Ben Mansour]{Jegham Mohamed \and Ben Mansour Louay}
\institute{\small Parcours Machine Learning}
\date{Avril 2026}

% ─────────────────── Intercalaire de section ─────────────────────────────────
\AtBeginSection[]{%
  \begin{frame}[plain]
    \begin{tikzpicture}[remember picture, overlay]
      % Fond dégradé simulé par deux rectangles
      \fill[MLDeep]    (current page.north west) rectangle (current page.south east);
      \fill[MLPrimary,opacity=0.45]
           ([xshift=6cm]current page.west) rectangle (current page.south east);
      % Numéro de section (grand, semi-transparent)
      \node[anchor=east, font=\fontsize{80}{80}\selectfont\bfseries,
            text=white, opacity=0.08]
           at ([xshift=-0.5cm]current page.east) {\insertsectionnumber};
    \end{tikzpicture}
    \vfill
    \centering
    \begin{tikzpicture}
      \node[%
        draw=white!40,%
        line width=1pt,%
        rounded corners=6pt,%
        inner sep=18pt,%
        text width=0.75\linewidth,%
        align=center,%
        fill=white!10,%
      ] {%
        \textcolor{white!60}{\small Section~\insertsectionnumber}\\[6pt]
        \textcolor{white}{\LARGE\bfseries\insertsectionhead}%
      };
    \end{tikzpicture}
    \vfill
  \end{frame}%
}

% =============================================================================
\begin{document}
% =============================================================================

% ── Page de titre ─────────────────────────────────────────────────────────────
\begin{frame}[plain]
  \begin{tikzpicture}[remember picture, overlay]
    % Bandeau supérieur
    \fill[MLDeep] (current page.north west)
         rectangle ([yshift=-1.6cm]current page.north east);
    % Ligne décorative accent
    \fill[MLAccent] ([yshift=-1.6cm]current page.north west)
         rectangle ([yshift=-1.76cm]current page.north east);
    % Date dans le bandeau
    \node[anchor=north east,
          font=\small\color{white!75}]
         at ([xshift=-0.8cm, yshift=-0.55cm]current page.north east)
         {Avril 2026};
    % Bande latérale gauche décorative
    \fill[MLPrimary!30] (current page.south west)
         rectangle ([xshift=0.28cm]current page.north west);
  \end{tikzpicture}

  \vspace*{1.5cm}
  \begin{minipage}[t]{0.92\linewidth}
    {\usebeamerfont{title}\usebeamercolor[fg]{title}\inserttitle\par}%
    \vspace{0.3em}
    \textcolor{MLPrimary!40}{\rule{0.18\linewidth}{1.2pt}}%
    \vspace{0.3em}\par
    {\usebeamerfont{subtitle}\usebeamercolor[fg]{subtitle}\insertsubtitle\par}
    \vspace{1.2em}
    {\usebeamerfont{author}\textbf{\insertauthor}\par}
    \vspace{0.3em}
    {\small\textcolor{MLMuted}{\insertinstitute}\par}
  \end{minipage}
\end{frame}

% ── Sommaire ──────────────────────────────────────────────────────────────────
\begin{frame}{Plan}
  \tableofcontents
\end{frame}

% =============================================================================
\section{Contexte et données}
% =============================================================================

\begin{frame}{Contexte}
  \begin{block}{Enjeux de la fraude bancaire}
    Volume élevé de transactions → nécessité d'un \textbf{scoring automatique} robuste.
  \end{block}
  \vspace{0.5em}
  \begin{itemize}
    \item Pertes financières, obligations réglementaires et confiance client en jeu.
    \item Classification binaire~: \textbf{légitime} \textit{vs} \textbf{frauduleux} sur données tabulaires réelles.
    \item Benchmark académique standard~: dataset ULB / Kaggle.
  \end{itemize}
\end{frame}

\begin{frame}{Jeu de données --- ULB / Kaggle}
  \begin{columns}[T,onlytextwidth]
    \column{0.46\textwidth}
      \begin{itemize}
        \item Fichier \texttt{creditcard.csv}~: $\approx 285\,000$ transactions.
        \item Variables~: \texttt{Time}, \texttt{Amount}, \texttt{V1}--\texttt{V28} (PCA anonymisée), cible \texttt{Class} $\in\{0,1\}$.
        \item \textbf{Déséquilibre sévère}~: $<0.2\,\%$ de fraudes.
        \item L'exactitude brute est \emph{trompeuse} → F1, ROC-AUC, Précision, Rappel.
      \end{itemize}
    \column{0.54\textwidth}
      \safeimg[0.97]{class_distribution.png}
  \end{columns}
\end{frame}

\begin{frame}{Exploration~: corrélations et distributions}
  \begin{columns}[T,onlytextwidth]
    \column{0.5\textwidth}
      \centering
      {\footnotesize\textcolor{MLMuted}{Corrélations (heatmap)}}\par\smallskip
      \safeimg[0.96]{correlation_heatmap.png}
    \column{0.5\textwidth}
      \centering
      {\footnotesize\textcolor{MLMuted}{Distributions des features}}\par\smallskip
      \safeimg[0.96]{feature_distributions.png}
  \end{columns}
\end{frame}

% =============================================================================
\section{Méthodologie}
% =============================================================================

\begin{frame}{Pipeline de traitement (vue d'ensemble)}
  \begin{enumerate}\setlength\itemsep{3pt}
    \item \textbf{EDA}~: distributions, corrélations, valeurs manquantes.
    \item \textbf{Split temporel}~: tri par \texttt{Time}, train = début, test = fin --- plus réaliste.
    \item \textbf{Features}~: \texttt{log1p(Amount)}, ratios vs moyenne mobile causale (sans fuite).
    \item \textbf{Sélection}~: régularisation L1 (logistique) sur le train.
    \item \textbf{Rééquilibrage}~: standardisation + \textbf{SMOTE} (train uniquement, \texttt{imblearn}).
    \item \textbf{Modèles}~: LR, Arbre, RF (GridSearch), XGBoost, LightGBM.
    \item \textbf{Seuil optimal}~: courbe PR + coût métier $C_\mathrm{FN}\cdot\mathrm{FN}+C_\mathrm{FP}\cdot\mathrm{FP}$.
    \item \textbf{Post-traitement}~: calibration, dérive (KS, PSI), SHAP.
    \item \textbf{Sauvegarde}~: \texttt{best\_model.pkl}, liste des variables \texttt{feature\_columns.json} (25~colonnes après sélection).
    \item \textbf{Démonstration}~: application \textbf{Streamlit} (seuil, graphiques Plotly, import CSV de test).
  \end{enumerate}
\end{frame}

\begin{frame}{SMOTE~: rééquilibrage sur le train}
  \centering
  \safeimg[0.72]{smote_class_distribution.png}
\end{frame}

% =============================================================================
\section{Résultats}
% =============================================================================

\begin{frame}{Comparaison des modèles (jeu de test)}
  {\small\textcolor{MLMuted}{Métriques issues de \texttt{outputs/model\_results.csv} (dernière exécution du notebook).}}
  \vspace{0.35em}
  \begin{center}
    \renewcommand{\arraystretch}{1.15}
    \scriptsize
    \begin{tabular}{@{}l *{5}{>{\centering\arraybackslash}p{1.05cm}}@{}}
      \toprule
      \rowcolor{MLDeep!90}
      \textcolor{white}{\textbf{Modèle}} &
      \textcolor{white}{\textbf{Acc.}} &
      \textcolor{white}{\textbf{Préc.}} &
      \textcolor{white}{\textbf{Rappel}} &
      \textcolor{white}{\textbf{F1}} &
      \textcolor{white}{\textbf{AUC}} \\
      \midrule
      Régression logistique & 0,984 & 0,065 & 0,853 & 0,121 & 0,981 \\
      \rowcolor{MLLight}
      Arbre de décision     & 0,993 & 0,132 & 0,773 & 0,225 & 0,820 \\
      \rowcolor{MLGreen!12}
      \textbf{Random Forest (tuné)} & \textbf{1,000} & \textbf{0,933} & \textbf{0,747} & \textbf{0,830} & \textbf{0,964} \\
      \rowcolor{MLLight}
      XGBoost               & 0,999 & 0,652 & 0,773 & 0,707 & 0,980 \\
      \rowcolor{MLLight}
      LightGBM              & 0,999 & 0,637 & 0,773 & 0,699 & 0,982 \\
      \bottomrule
    \end{tabular}
  \end{center}
  \vspace{0.4em}
  \begin{exampleblock}{Lecture}
    \small Meilleur \textbf{F1} sur ce run~: \textbf{Random Forest} (ligne verte). Les valeurs se mettent à jour automatiquement après \texttt{02\_Modeling.ipynb}.
  \end{exampleblock}
\end{frame}

\begin{frame}{Matrice de confusion~: modèle retenu}
  \centering
  \safeimg[0.62]{confusion_matrix.png}
\end{frame}

\begin{frame}{Courbe ROC}
  \centering
  \safeimg[0.78]{roc_curve.png}
\end{frame}

\begin{frame}{Précision–Rappel et choix de seuil}
  \centering
  \safeimg[0.78]{precision_recall_curve.png}
\end{frame}

\begin{frame}{Calibration des probabilités}
  \centering
  \safeimg[0.78]{calibration_curve.png}
\end{frame}

\begin{frame}{Importance des variables (modèle retenu)}
  \centering
  \safeimg[0.78]{feature_importance.png}
\end{frame}

\begin{frame}{Explicabilité~: SHAP (summary plot)}
  \centering
  \safeimg[0.78]{shap_summary.png}
\end{frame}

% =============================================================================
\section{Sorties du projet et démonstration}
% =============================================================================

\begin{frame}[t]{Artefacts dans \texttt{outputs/}}
  \begin{columns}[T,onlytextwidth]
    \column{0.52\textwidth}
      \begin{block}{Fichiers générés par le pipeline}
        \small
        \begin{itemize}
          \item \textbf{Modèle}~: \texttt{best\_model.pkl}
          \item \textbf{Entrées du scoreur}~: \texttt{feature\_columns.json} (25~variables retenues)
          \item \textbf{Tableau de synthèse}~: \texttt{model\_results.csv}
          \item \textbf{Dérive}~: \texttt{drift\_report.csv}
          \item \textbf{Figures}~: ROC, PR, calibration, SHAP, matrices\ldots
        \end{itemize}
      \end{block}
    \column{0.48\textwidth}
      \begin{exampleblock}{Cohérence}
        \small L’interface Streamlit charge le même pipeline et les mêmes noms de colonnes que le notebook — pas de décalage de schéma.
      \end{exampleblock}
      \vspace{0.3em}
      {\footnotesize\textcolor{MLMuted}{Relancer \texttt{02\_Modeling.ipynb} après toute modification du flux pour régénérer ces fichiers.}}
  \end{columns}
\end{frame}

\begin{frame}{Dérive train \textit{vs} test (\texttt{drift\_report.csv})}
  \small
  Extrait des indicateurs exportés (KS, PSI) — \textit{valeurs typiques d’une exécution récente}~:
  \vspace{0.4em}
  \begin{center}
    \renewcommand{\arraystretch}{1.15}
    \begin{tabular}{@{}lccc@{}}
      \toprule
      \rowcolor{MLDeep!85}
      \textcolor{white}{\textbf{Variable}} &
      \textcolor{white}{\textbf{KS}} &
      \textcolor{white}{\textbf{p-val. KS}} &
      \textcolor{white}{\textbf{PSI}} \\
      \midrule
      \texttt{Amount\_log} & 0,053 & $\approx 0$ & 0,014 \\
      \texttt{Time}        & 1,000 & $\approx 0$ & 11,51 \\
      \texttt{V14}         & 0,086 & $\approx 0$ & 0,061 \\
      \texttt{V12}         & 0,088 & $\approx 0$ & 0,151 \\
      \bottomrule
    \end{tabular}
  \end{center}
  \vspace{0.35em}
  \begin{alertblock}{Lecture}
    \small Split temporel~: \texttt{Time} peut fortement diverger entre train et test (PSI élevé). Sur les \texttt{V*} et \texttt{Amount\_log}, le PSI reste souvent modéré — à suivre en production (seuil courant PSI $>0{,}2$).
  \end{alertblock}
\end{frame}

\begin{frame}{Interface Streamlit — essai du modèle}
  \begin{itemize}
    \item \textbf{Seuil} réglable + coût métier $C_{\mathrm{FN}}{:}C_{\mathrm{FP}}$ aligné sur \texttt{src/config.py}.
    \item \textbf{Échantillon}~: fenêtre aléatoire, depuis le début, ou décalage fixe — reproductibilité affichée, export CSV des prédictions.
    \item \textbf{Visualisation}~: histogrammes Plotly des probabilités par vraie classe.
    \item \textbf{Import CSV}~: fichier avec \texttt{Time}, \texttt{Amount}, \texttt{V1}--\texttt{V28} ; historique causal du jeu pour les moyennes mobiles ; résultats téléchargeables.
  \end{itemize}
  \vspace{0.5em}
  \begin{block}{Lancement}
    \centering\texttt{streamlit run streamlit\_app.py}\quad\textit{(depuis la racine du projet)}
  \end{block}
\end{frame}

% =============================================================================
\section{Limites et perspectives}
% =============================================================================

\begin{frame}{Limites et perspectives}
  \begin{columns}[T,onlytextwidth]
    \column{0.5\textwidth}
      \begin{alertblock}{Limites}
        \begin{itemize}
          \item Données de 2013~: risque de dérive en production.
          \item Variables PCA anonymisées~: interprétation métier restreinte.
        \end{itemize}
      \end{alertblock}
    \column{0.5\textwidth}
      \begin{exampleblock}{Pistes}
        \begin{itemize}
          \item CatBoost, TabNet.
          \item Coûts métier personnalisés.
          \item MLOps~: monitoring PSI/KS continu.
        \end{itemize}
      \end{exampleblock}
  \end{columns}
\end{frame}

% =============================================================================
\section{Conclusion}
% =============================================================================

\begin{frame}{Conclusion}
  \begin{block}{Livrables techniques}
    \small
    \texttt{requirements.txt} · \texttt{01\_EDA} / \texttt{02\_Modeling} · \texttt{src/} · \texttt{outputs/*} · \texttt{streamlit\_app.py} · tests \texttt{pytest}
  \end{block}
  \vspace{0.5em}
  \begin{itemize}
    \item Automatisation~: \texttt{Makefile}, \texttt{tasks.ps1}~; CI GitHub Actions.
    \item \textbf{Random Forest (GridSearch)}~: meilleur F1 sur le jeu de test actuel (\textbf{F1 $\approx 0{,}83$}, AUC-ROC $\approx 0{,}96$) — voir \texttt{model\_results.csv}.
    \item \textbf{Démonstration interactive}~: scoring, seuil, import CSV et figures Plotly via Streamlit.
    \item Le modèle reste un \textbf{outil d'aide à la décision}~: seuil et arbitrage FN/FP sont métier.
  \end{itemize}
\end{frame}

% ── Diapo de fin ──────────────────────────────────────────────────────────────
\begin{frame}[plain]
  \begin{tikzpicture}[remember picture, overlay]
    \fill[MLDeep]    (current page.south west) rectangle (current page.north east);
    \fill[MLPrimary,opacity=0.4]
         ([xshift=5cm]current page.west) rectangle (current page.south east);
    \fill[MLAccent]  (current page.south west)
         rectangle ([yshift=0.18cm]current page.south east);
  \end{tikzpicture}
  \vspace*{2.0cm}
  \centering
  {\LARGE\bfseries\color{white}Merci de votre attention}\par
  \vspace{1.2em}
  \textcolor{white!50}{\rule{0.20\linewidth}{0.6pt}}\par
  \vspace{1.0em}
  \href{https://www.kaggle.com/datasets/mlg-ulb/creditcardfraud}{%
    \textcolor{MLAccent!70}{\small Dataset Kaggle --- ULB Credit Card Fraud}%
  }\par
  \vspace{0.6em}
  {\small\textcolor{white!60}{\insertauthor\ $\cdot$\ \insertinstitute}}
\end{frame}

\end{document}